%-----------------------------------------------------------------------------------------------------------------------------------------------%
%	The MIT License (MIT)
%
%	Copyright (c) 2021 Jitin Nair
%
%	Permission is hereby granted, free of charge, to any person obtaining a copy
%	of this software and associated documentation files (the "Software"), to deal
%	in the Software without restriction, including without limitation the rights
%	to use, copy, modify, merge, publish, distribute, sublicense, and/or sell
%	copies of the Software, and to permit persons to whom the Software is
%	furnished to do so, subject to the following conditions:
%	
%	THE SOFTWARE IS PROVIDED "AS IS", WITHOUT WARRANTY OF ANY KIND, EXPRESS OR
%	IMPLIED, INCLUDING BUT NOT LIMITED TO THE WARRANTIES OF MERCHANTABILITY,
%	FITNESS FOR A PARTICULAR PURPOSE AND NONINFRINGEMENT. IN NO EVENT SHALL THE
%	AUTHORS OR COPYRIGHT HOLDERS BE LIABLE FOR ANY CLAIM, DAMAGES OR OTHER
%	LIABILITY, WHETHER IN AN ACTION OF CONTRACT, TORT OR OTHERWISE, ARISING FROM,
%	OUT OF OR IN CONNECTION WITH THE SOFTWARE OR THE USE OR OTHER DEALINGS IN
%	THE SOFTWARE.
%	
%
%-----------------------------------------------------------------------------------------------------------------------------------------------%

%----------------------------------------------------------------------------------------
%	DOCUMENT DEFINITION
%----------------------------------------------------------------------------------------

% article class because we want to fully customize the page and not use a cv template
\documentclass[a4paper,12pt]{article}

%----------------------------------------------------------------------------------------
%	FONT
%----------------------------------------------------------------------------------------

% % fontspec allows you to use TTF/OTF fonts directly
% \usepackage{fontspec}
% \defaultfontfeatures{Ligatures=TeX}

% % modified for ShareLaTeX use
% \setmainfont[
% SmallCapsFont = Fontin-SmallCaps.otf,
% BoldFont = Fontin-Bold.otf,
% ItalicFont = Fontin-Italic.otf
% ]
% {Fontin.otf}

%----------------------------------------------------------------------------------------
%	PACKAGES
%----------------------------------------------------------------------------------------
\usepackage{url}
\usepackage{parskip} 	

%other packages for formatting
\RequirePackage{color}
\RequirePackage{graphicx}
\usepackage[usenames,dvipsnames]{xcolor}
\usepackage[scale=0.9]{geometry}

%tabularx environment
\usepackage{tabularx}

%for lists within experience section
\usepackage{enumitem}

% centered version of 'X' col. type
\newcolumntype{C}{>{\centering\arraybackslash}X} 

%to prevent spillover of tabular into next pages
\usepackage{supertabular}
\usepackage{tabularx}
\newlength{\fullcollw}
\setlength{\fullcollw}{0.47\textwidth}

%custom \section
\usepackage{titlesec}				
\usepackage{multicol}
\usepackage{multirow}

% Pour que la date s'affiche en Français : 
\usepackage[francais]{babel}

%CV Sections inspired by: 
%http://stefano.italians.nl/archives/26
\titleformat{\section}{\Large\scshape\raggedright}{}{0em}{}[\titlerule]
\titlespacing{\section}{0pt}{10pt}{10pt}

%for publications
\usepackage[style=authoryear,sorting=ynt, maxbibnames=2]{biblatex}

%Setup hyperref package, and colours for links
\usepackage[unicode, draft=false]{hyperref}
\definecolor{linkcolour}{rgb}{0,0.2,0.6}
\hypersetup{colorlinks,breaklinks,urlcolor=linkcolour,linkcolor=linkcolour}
\addbibresource{citations.bib}
\setlength\bibitemsep{1em}

%for social icons
\usepackage{fontawesome5}

%debug page outer frames
%\usepackage{showframe}

%----------------------------------------------------------------------------------------
%	BEGIN DOCUMENT
%----------------------------------------------------------------------------------------
\begin{document}

% non-numbered pages
\pagestyle{empty} 

%----------------------------------------------------------------------------------------
%	TITLE
%----------------------------------------------------------------------------------------

% \begin{tabularx}{\linewidth}{ @{}X X@{} }
% \huge{Your Name}\vspace{2pt} & \hfill \emoji{incoming-envelope} email@email.com \\
% \raisebox{-0.05\height}\faGithub\ username \ | \
% \raisebox{-0.00\height}\faLinkedin\ username \ | \ \raisebox{-0.05\height}\faGlobe \ mysite.com  & \hfill \emoji{calling} number
% \end{tabularx}

\begin{tabularx}{\linewidth}{@{} C @{}}
\Huge{Paulo Gugelmo Cavalheiro Dias} \\[7.5pt]
\href{https://github.com/Paulogcd}{\raisebox{-0.05\height}\faGithub\ Paulogcd} \ $|$ \ 
\href{https://www.linkedin.com/in/paulo-gugelmo-cavalheiro-dias/}{\raisebox{-0.05\height}\faLinkedin\ LinkedIn} \ $|$ \ 
\href{https://paulogcd.com}{\raisebox{-0.05\height}\faGlobe \ paulogcd.com} \ $|$ \ 
\href{mailto:paulo.gugelmocavalheirodias@sciencespo.fr}{\raisebox{-0.05\height}\faEnvelope \ Email} \ 
\end{tabularx}

%----------------------------------------------------------------------------------------
% EXPERIENCE SECTIONS
%----------------------------------------------------------------------------------------

\hfill \break

%Experience
\section{Expériences professionnelles}

\begin{tabularx}{\linewidth}{ @{}l r@{} }
\textbf{Global Impact Metrics - Consultant} & \hfill Novembre 2024 - Présent \\
\multicolumn{2}{@{}X@{}}{
    \begin{itemize}
    \item Analyse quantitative et gestion de bases de données de grande dimension.
    \item Analyse statistique et mise en œuvre de flux de travail reproductibles.
    \item Revues de littérature, synthèse des connaissances et analyse des politiques publiques.
    \item Évaluation d'impact et analyses coûts-bénéfices.
    \item Gestion et déploiement de services web internes.
    \end{itemize}
} \
\end{tabularx}

\begin{tabularx}{\linewidth}{ @{}l r@{} }
\textbf{Centre de recherche Marc Bloch} & \hfill Avril 2023 - Août 2023 \\[3.75pt]
\multicolumn{2}{@{}X@{}}{
    \begin{itemize}
        \item Assistance à la recherche et revues de littérature.
        \item Coordination d'événements académiques et soutien administratif.
    \end{itemize}
} \
\end{tabularx}

\begin{tabularx}{\linewidth}{ @{}l r@{} }
\textbf{Chambre Franco-Allemande de Commerce et d'Industrie} & \hfill Septembre 2022 - Mars 2023 \\[3.75pt]
\multicolumn{2}{@{}X@{}}{
\begin{itemize}
    \item Analyse économique des entreprises.
    \item Collecte, gestion et reporting de données commerciales.
\end{itemize}
}
\end{tabularx}

\begin{tabularx}{\linewidth}{ @{}l r@{} }
\textbf{Mairie de Beauchamp} & \hfill Juillet 2022 – Septembre 2022 \\[3.75pt]
\multicolumn{2}{@{}X@{}}{
\begin{itemize}
    \item Analyse économique et évaluation des politiques publiques locales.
    \item Collecte de données et production de rapports statistiques.
\end{itemize}    
}  \\
\end{tabularx}

\begin{tabularx}{\linewidth}{ @{}l r@{} }
\textbf{Sciences Po, Centre de Sociologie des Organisations (CSO)} & \hfill Juillet 2021 – Septembre 2021 \\[3.75pt]
\multicolumn{2}{@{}X@{}}{
\begin{itemize}
    \item Collecte de données via l'API Twitter avec R.
    \item Traitement, gestion et préparation des données pour l'analyse empirique.
    \item Visualisation de réseaux à l'aide du package igraph.
\end{itemize}
}  \\
\end{tabularx}

%----------------------------------------------------------------------------------------
%	EDUCATION
%----------------------------------------------------------------------------------------
\section{Formation}
\begin{tabularx}{\linewidth}{@{}l X@{}}	

2023 – 2025 & \textit{Master of Research in Economics} à \href{https://www.sciencespo.fr/en/}{Sciences Po} : \\ 
& \begin{itemize}
    \item Macroéconomie, microéconomie, économétrie - niveaux I, II, et III,
    \item Programmation : gestion de projet, R, Julia, Python,
    \item Économie financière, du travail, publique,
    \item Commerce international, économie politique, du développement, et organisationnelle.
\end{itemize} \\

2019 – 2023 & Double \textit{Bachelor} en sciences sociales et science politique à \href{https://www.sciencespo.fr/en/}{Sciences Po} et à la \textit{\href{https://www.fu-berlin.de/en/index.html}{Freie Universität}} (mention Bien)\\

2019 – 2022 & Licence en économie et gestion (mention Bien) à l'\href{https://www.univ-lorraine.fr/en/}{Université de Lorraine}\\

2019 & Baccalauréat au \href{https://www.lyc-polyvalent-monod-enghien.fr/}{lycée Gustave Monod, Enghien-les-Bains} \hfill \\

2019 & Abitur au \href{https://www.lyc-polyvalent-monod-enghien.fr/}{lycée Gustave Monod, Enghien-les-Bains} \hfill \\
\end{tabularx}

%----------------------------------------------------------------------------------------
%	SKILLS
%----------------------------------------------------------------------------------------
\section{Compétences}
\begin{tabularx}{\linewidth}{@{}l X@{}}
Programmation 
& Niveau avancé avec R en analyse de données :
\begin{itemize}
\item Collecte de données (rvest, httr),
\item Analyse statistique (packages tydiverse, data.table, igraph pour l'analyse de réseau),
\item Visualisation de données (ggplot2, rmarkdown, notions en Shiny).
\end{itemize} \\
& Niveau avancé avec Julia en traitement de données et modélisation économique : 
\begin{itemize}
\item Analyse statistique (calculs, production de table de régressions, etc...),
\item Modélisation économique (modèle macroéconomique, simulations, résolution numérique d'équations transcendantales), 
\item Intégration des résultats via Plotly et Quarto,
\item Méthodes numériques (value-function iteration, policy-function iteration, évaluation algorithmique, méthodes de parallélisme e.g. multi-threading).
\end{itemize}\\
& Niveau élémentaire avec d'autres logiciels statistiques : Stata, MATLAB. \\
& Niveau élémentaire en Python, SQL, et C. \\
& \\
Développement web &  
\begin{itemize}
    \item Gestion et déploiement de services et applications web avec Docker et Nginx.
    \item Intégration et déploiement continus (GitLab CI/CD, GitHub Actions).
    \item Expérience du framework Quarto (voir mon \href{https://www.paulogcd.com}{site internet}).
\end{itemize} \\ 
Autres & Maîtrise de la suite Microsoft Office. \\
&  \LaTeX. \\
& Connaissances basiques en shell et de l'architecture Unix. \\
& Gestion de code et contrôle de version avec git.\\
\end{tabularx}

\section{Projets sélectionnés}

\begin{tabularx}{\linewidth}{@{}l X@{}}
    \href{https://www.paulogcd.com/Master_Thesis}{Site web et package de réplication du mémoire de master} &
    Site compagnon développé autour d'un dépôt Git (et déployé avec GitHub pages) pour un mémoire de master à dominante computationnelle.
    Réalisé avec Julia, Pluto.jl, Quarto et Plotly, il fournit des visualisations interactives, des notebooks reproductibles ainsi qu'un package complet de réplication des résultats. \\ \\

    \href{https://www.paulogcd.com/Replication_Monge_et_al_2019.jl/dev/}{Réplication de Monge-Naranjo et al. (2019)} &
    Projet collectif consistant à traduire en Julia le package de réplication Stata de l'article Natural Resources and Global Misallocation (American Economic Journal: Macroeconomics, 2019), dans le cadre du cours de Computational Economics de Florian Oswald à Sciences Po.
    Le projet met l'accent sur la reproductibilité et l'utilisation de Julia en économie computationnelle. \\ \\

    \href{https://www.paulogcd.com/portfolio_tutorial/}{Tutoriel de portfolio de recherche} &
    Tutoriel (en anglais) présentant la création et la maintenance d'un portfolio de recherche en ligne à l'aide de VS Code, Git, GitHub, Quarto et GitHub Pages.
    Conçu pour aider les étudiants en recherche à mettre en place des flux de travail reproductibles et à valoriser leurs travaux. \\ \\

\end{tabularx}

%----------------------------------------------------------------------------------------
%	LANGUAGES
%----------------------------------------------------------------------------------------
\section{Langues}
\begin{tabularx}{\linewidth}{@{}l X@{}}
Français &  \normalsize{Natif (C2)}\\
Portugais &  \normalsize{Parlé (C2)}\\
Allemand &  \normalsize{Niveau professionnel (C1)}\\  
Anglais &  \normalsize{Niveau professionnel (C1) - TOEFL : 104/120}\\ 
Espagnol & \normalsize{Niveau débutant (B1)}\\ 
Italien & \normalsize{Niveau débutant (B1)}\\ 
\end{tabularx}

%----------------------------------------------------------------------------------------
%	Références
%----------------------------------------------------------------------------------------
\section{Références}

\href{https://floswald.github.io}{Prof. Florian Oswald} - florian.oswald@gmail.com \\
\href{https://www.jcommault.com}{Prof. Jeanne Commault} - jeanne.commault@sciencespo.fr \\
\href{https://www.paul-bouscasse.com}{Prof. Paul Bouscasse} - paul.bouscasse@sciencespo.fr \\

\vfill
\center{\footnotesize Dernière mise à jour: \selectlanguage{french}\today}

\end{document}